
\section{Introduction}
\label{sec:intro}

Solar generation is scheduled against forward-looking forecasts of global horizontal irradiance (GHI) and co-dependent boundary-layer meteorological variables. When deterministic forecasts diverge from atmospheric reality, grid operators must balance real-time deficits or surpluses using expensive operating reserves, localized generation curtailment, or fast-ramping fossil-fuel resources. These balancing actions incur substantial economic penalties, strain regional transmission infrastructure, and undermine the carbon-mitigation benefits of renewable integration. 

For power system planning, market clearing, and reserve dimensioning, the relevant metric is therefore not merely global forecast accuracy, but a structural characterization of the underlying error fields. Specifically, operators must understand how forecast error variance expands across the operational timeline, whether errors propagate with short-term temporal persistence, and if spatially distributed assets offer error diversification—where independent, co-located prediction misses cancel out under portfolio aggregation rather than reinforcing one another.

This report establishes a rigorous empirical evaluation of the High-Resolution Rapid Refresh (HRRR) model's error structures across the 896-site PJM solar fleet using a pooled 2022–2024 dataset. We systematically investigate the boundary conditions of error predictability by addressing four foundational questions:
\begin{enumerate}
  \item \textbf{How does forecast error accumulate, distribute, and stabilize across the 48-hour horizon?} We examine the baseline error distributions using robust Kernel Density Estimation (KDE) and implement a parametric sequential Cumulative Sum (CUSUM) change-point framework to isolate the mathematical boundaries where early-stage skill degradation transitions into a stabilized uncertainty plateau (Section~\ref{sec:descriptive}).
  \item \textbf{What drives hourly error growth rates, and are they systematically correctable?} We evaluate fleet-mean error growth trends to determine if forecast horizons introduce monotonic bias drift or if they are dominated by cyclic, diurnal confounding factors that can be modeled via harmonic regression (Section~\ref{sec:leadtime}).
  \item \textbf{What is the internal temporal memory structure of these errors?} Using the Autocorrelation Function (ACF) and Partial Autocorrelation Function (PACF) combined with non-parametric Kruskal-Wallis testing, we quantify the short-term persistence of errors and assess whether this temporal structure is stable across seasons, years, and transmission zones.
  \item \textbf{How do forecast errors correlate and decay across space, and does this coherence vary by environmental stratum?} We map zone-to-zone Pearson correlations and fit exponential decay models to determine the characteristic decorrelation lengths ($L$) per variable, establishing how spatial independence shifts across hours of the day, seasonal weather regimes, and subsequent calendar years (Section~\ref{sec:spatial} and Section~\ref{sec:strata}).
\end{enumerate}

The structural thread unifying this investigation is the mathematical concept of \emph{portfolio diversification}. Meteorological forecast errors only average out across a regional transmission footprint when the errors at constituent generation nodes approach statistical independence. By identifying the explicit temporal inflections, persistence boundaries, and spatial e-folding decay scales unique to each atmospheric variable, this report delivers the empirical foundation required to build advanced bias-correction filters and portfolio-level risk-aggregation models.